\documentclass{TesiDiaUniroma3}
\usepackage[utf8]{inputenc}
\usepackage{graphicx}
\usepackage[bottom]{footmisc}

\titolo{\textbf{Elettroencefalografia per interfacce cervello-computer: classificazione dell'immaginazione motoria e predizione delle prestazioni}}

\autore{Alessio Paparcurio}
\matricola{634690}
\relatore{Prof. Fabio Gasparetti}
\correlatore{Prof.ssa Silvia Conforto}
\annoAccademico{2025/2026}


% --- FINE dati relativi al template TesiDiaUniroma3

% --- INIZIO richiamo di pacchetti di utilità. Questi sono un esempio e non sono strettamente necessari al modello per la tesi.
\usepackage[plainpages=false]{hyperref}	% generazione di collegamenti ipertestuali su indice e riferimenti
\usepackage[all]{hypcap} % per far si che i link ipertestuali alle immagini puntino all'inizio delle stesse e non alla didascalia sottostante
\usepackage{amsthm}	% per definizioni e teoremi
\usepackage{amsmath}	% per ``cases'' environment
\usepackage{graphicx}
\usepackage{float}
\usepackage{hyperref}
\usepackage{amsfonts}
\usepackage{subcaption}
\usepackage{mathtools}
\usepackage[width=.75\textwidth]{caption}
\usepackage{adjustbox}
\usepackage[normalem]{ulem}
\usepackage[]{algorithm2e}

% Package aggiunti 
\usepackage{makecell} % Per le tabelle
\usepackage{listings}
\usepackage[table]{xcolor}
\usepackage{rotating}
\usepackage{tabularx}
\usepackage{multicol}
\usepackage{booktabs}
\usepackage{lscape}
\usepackage[flushleft]{threeparttable}
\usepackage[autostyle]{csquotes}

\useunder{\uline}{\ul}{}
% Impostazioni per i frammenti di codice.
\lstset{
	language=Python,
	basicstyle=\ttfamily\footnotesize\SingleSpacing,
	keywordstyle=\bfseries,
	commentstyle=\itshape,
	showstringspaces=false,
	breaklines=true,
	frame=single,
	xleftmargin=1em,
	aboveskip=1.2em,
	belowskip=1.2em
}
% --- FINE richiamo di pacchetti di utilità

\setlength{\emergencystretch}{3em} % evita righe che sbordano nel margine

\begin{document}
% ----- Pagine di frontespizio, numerate in romano (i,ii,iii,iv...) (obbligatorio)
\frontmatter
\generaFrontespizio

% \capitolo{Abstract}{capitoli/abstract}
\generaIndice
\generaIndiceFigure
\generaIndiceTabelle

% ----- Pagine di tesi, numerate in arabo (1,2,3,4,...) (obbligatorio)
\mainmatter

\introduzione{capitoli/01-introduzione}
\capitolo{Principali tecnologie e paradigmi moderni}{capitoli/02-stato-arte}
\capitolo{Dataset}{capitoli/03-dataset}
\capitolo{Classificatore e predittore di illiteracy}{capitoli/04-pipeline}
\capitolo{Risultati sperimentali}{capitoli/05-risultati}
\capitolo{Conclusioni e sviluppi futuri}{capitoli/06-conclusioni}

% ----- Parte finale della tesi (obbligatorio)
\backmatter
\appendix
% \capitolo{Appendice}{capitoli/appendice}
% \ringraziamenti{capitoli/ringraziamenti}
\bibliography{bibliografia} 

\end{document}
